% ---------------------------------------------------------------------------
% paper/drafts/result_transition_paragraphs.tex
% Short connective paragraphs between the results sections.
% Draft fragment. NOT included in the main paper.
% Each block is independently usable; keep only the ones the final section
% ordering needs. NONE of these asserts a pending outcome.
% ---------------------------------------------------------------------------

% --- Method -> matched-synthetic results -----------------------------------
\newcommand{\TransitionMethodToMatched}{%
Everything so far has been an argument about what the model and its sampler are
supposed to do. The matched-synthetic study asks whether they do it, under a
generator that implements the inference target exactly rather than
approximately: the same normalised segmentation prior, the same per-block
recurrent reset, the same zero-diagonal transition law. That match is itself a
validated artifact, so a recovery failure on this corpus would be a statement
about identifiability rather than about a generator/inference mismatch.%
}

% --- Condition A -> Condition B --------------------------------------------
\newcommand{\TransitionAToB}{%
Condition~A establishes path identifiability under oracle structures: with the
reusable structures and global parameters held at truth, the exact posterior
over skill-instance boundaries and labels is sharply informative and well
calibrated. Condition~B inverts the conditioning and establishes
reusable-structure identifiability under oracle paths: with the skill-instance
segmentation and labels held at truth, four chains started from four distinct
wrong structures all concentrate on the generating structure.%
}

% --- Conditions A and B -> Condition C -------------------------------------
\newcommand{\TransitionABToC}{%
Taken together, the two conditions isolate what is left to test. Neither
component is unidentifiable on its own, so the open question is not whether the
data carry information about paths or about structures, but whether the joint
posterior over both can be explored when neither is observed --- a question
about posterior coupling and sampler geometry rather than about component
identifiability. Condition~C runs exactly that experiment, with both $(S,z)$ and
$H$ latent, in two arms that target the same posterior: a conditional arm and a
path-marginal arm.%
}

% --- Condition C -> external ------------------------------------------------
\newcommand{\TransitionCToExternal}{%
The matched study evaluates HPOP inference under controlled truth, where every
latent quantity has a known answer and the generator is the model. That is the
right setting for identifiability and for sampler validation, and the wrong
setting for generalisation: a model can be perfectly recoverable on data drawn
from itself and still describe nothing outside it. The external experiments
therefore ask a different question --- whether the representation transfers to
data the model did not generate.%
}

% --- Matched -> TaskBench (usable if Condition C is not yet in the paper) ---
\newcommand{\TransitionMatchedToTaskbench}{%
The matched study evaluates HPOP inference under controlled truth; TaskBench
evaluates whether the static partial-order representation generalises beyond
internally generated data. The two answer different questions and neither
substitutes for the other: TaskBench supplies real external task structure but
no latent segmentation, so it can test the partial-order component and nothing
above it.%
}

% --- TaskBench -> tau^3-Retail ---------------------------------------------
\newcommand{\TransitionTaskbenchToTauThree}{%
TaskBench validates only the static component, because its orderings are
enumerated from gold graphs rather than observed and it contains no repeated
skill instances or repair behaviour. $\tau^3$-Retail provides the trajectory
structure that a full-hierarchy external pilot needs: real agent tool sequences,
a repeated-CPA fraction above $0.9$, both successful and failed runs, and a
task-disjoint split. It provides no gold segmentation or gold partial order, so
it is evaluated predictively rather than by structural recovery.%
}

% --- Scope reminder, usable as a standalone limitations lead-in -------------
\newcommand{\TransitionScopeReminder}{%
Two of our experimental settings speak to HPOP as a whole and two speak to parts
of it. The matched-synthetic conditions test the full generative model under
controlled truth; the collapsed-$U$ validation tests the sampler that makes
joint inference tractable. TaskBench tests only the static partial-order
component on external data, and the $\tau^3$-Retail package is at present a
validated corpus rather than a result. We keep those scopes separate throughout,
and Table~\ref{tab:claim-scope} states for each experiment what is established
and what is not.%
}

% ---------------------------------------------------------------------------
% NOTE TO AUTHOR (delete before submission):
%  - \TransitionABToC describes what Condition C ASKS. It must not be edited
%    into a claim about what Condition C FOUND until a frozen artifact exists.
%  - \TransitionTaskbenchToTauThree describes the corpus only. Do not append a
%    predictive claim to it.
% ---------------------------------------------------------------------------
